
\documentclass{article}
\usepackage{colortbl}
\usepackage{makecell}
\usepackage{multirow}
\usepackage{supertabular}

\begin{document}

\newcounter{utterance}

\centering \large Interaction Transcript for game `codenames', experiment `abstract', episode 4 with reflexion{-}Qwen3.5{-}9B.
\vspace{24pt}

{ \footnotesize  \setcounter{utterance}{1}
\setlength{\tabcolsep}{0pt}
\begin{supertabular}{c@{$\;$}|p{.15\linewidth}@{}p{.15\linewidth}p{.15\linewidth}p{.15\linewidth}p{.15\linewidth}p{.15\linewidth}}
    \# & $\;$A & \multicolumn{4}{c}{Game Master} & $\;\:$B\\
    \hline

    \theutterance \stepcounter{utterance}  
    & & & \multicolumn{2}{p{0.3\linewidth}}{
        \cellcolor[rgb]{0.9,0.9,0.9}{
            \makecell[{{p{\linewidth}}}]{
                \texttt{\tiny{[GM$|$GM]}}
                {'hidden': {'team': ['bonkers', 'ambivalence'], 'innocent': ['nationalization', 'impulsivity', 'socialization'], 'opponent': [], 'assassin': []}, 'revealed': {'team': {'team': [], 'innocent': [], 'opponent': [], 'assassin': []}, 'opponent': {'team': [], 'innocent': [], 'opponent': [], 'assassin': []}}}
            }
        }
    }
    & & \\ \\

    \theutterance \stepcounter{utterance}  
    & & \multicolumn{4}{p{0.6\linewidth}}{
        \cellcolor[rgb]{0.9,0.9,0.9}{
            \makecell[{{p{\linewidth}}}]{
                \texttt{\tiny{[P1$\langle$GM]}}
                \texttt{You have attempted this same task before and did not succeed. These are the plans you wrote after those attempts.} \\
\texttt{Trial \#1: New plan: I failed because my clues were too broad, causing my partner to guess distractor words (like "gladness" and "subjugation") which ended their turn prematurely. I need to select clue words that are semantically specific to the remaining target words and have no overlap with the distractor list. For the next turn, I will choose a highly specific clue like "chaos" for "bonkers" or "outcome" for "fruition" to ensure the partner guesses the correct target immediately without risking an innocent word. I will also limit the number of targets per clue to reduce ambiguity.} \\
\\ 
\texttt{Trial \#2: New plan: I failed because my clues were too broad (e.g., "chaos", "mixed"), causing my partner to guess innocent distractor words like "proliferation" and "incompatibility," which prematurely ended my turn. I need to select clue words that are semantically specific to the remaining target words and have no overlap with the distractor list. For the next turn, I will choose a highly specific clue like "crazy" for "bonkers" or "feeling" for "ambivalence" to ensure the partner guesses the correct target immediately without risking an innocent word. I will also limit the number of targets per clue to reduce ambiguity.} \\
\\ 
\texttt{Trial \#3: New plan: I failed because my clue "mental" was too broad, allowing my partner to guess distractor words like "impulsivity" or "socialization," which prematurely ended the turn. Next time, I will choose a highly specific clue that directly defines the target words without overlapping with distractors. For "bonkers," I will use "crazy" or "unhinged," and for "ambivalence," I will use "mixed feelings" or "contradiction." I will also limit the number of targets per clue to one to ensure clarity and reduce the risk of guessing errors.} \\
\\ 
\\ 
\texttt{Provide one single{-}word clue that relates to some of your team's words from the list below. You can choose to how many and to which words your clue relates to.} \\
\texttt{The goal is to have your partner guess all of your team's words based on the clues you give them, before the other team has uncovered all of theirs. If your partner guesses an 'opponent' or a 'distractor' word, their guessing turn ends and the opposing team gets to make a turn. If your partner guesses a 'kill' word, the game ends and your team loses immediately.} \\
\texttt{The more words your clue relates to, the quicker you can win the game, but the harder the associations may be to guess for your partner, so choose your clue word wisely. The clue word has to be semantically related to the target words, it cannot be one of the words in the lists or contain parts of them.} \\
\\ 
\texttt{Always give your single{-}word clue and your comma{-}separated list of related target words in the following format and make your answers as short as possible, never include any other text than is required in this form:} \\
\\ 
\texttt{CLUE: <WORD>} \\
\texttt{TARGETS: <TARGETS>} \\
\\ 
\texttt{Your team words are: bonkers, ambivalence.} \\
\texttt{Your opponent's team words are: .} \\
\texttt{Distractor words are: nationalization, impulsivity, socialization.} \\
\texttt{The kill words are: .} \\
            }
        }
    }
    & \\ \\

    \theutterance \stepcounter{utterance}  
    & \multicolumn{4}{p{0.6\linewidth}}{
        \cellcolor[rgb]{0.9,0.9,0.9}{
            \makecell[{{p{\linewidth}}}]{
                \texttt{\tiny{[P1$\rangle$GM]}}
                \texttt{CLUE: contradiction} \\
\texttt{TARGETS: ambivalence, bonkers} \\
            }
        }
    }
    & & \\ \\

    \theutterance \stepcounter{utterance}  
    & & & \multicolumn{2}{p{0.3\linewidth}}{
        \cellcolor[rgb]{0.9,0.9,0.9}{
            \makecell[{{p{\linewidth}}}]{
                \texttt{\tiny{[GM$|$GM]}}
                \texttt{contradiction} \\
            }
        }
    }
    & & \\ \\

    \theutterance \stepcounter{utterance}  
    & & & \multicolumn{2}{p{0.3\linewidth}}{
        \cellcolor[rgb]{0.9,0.9,0.9}{
            \makecell[{{p{\linewidth}}}]{
                \texttt{\tiny{[GM$|$GM]}}
                ['ambivalence', 'bonkers']
            }
        }
    }
    & & \\ \\

    \theutterance \stepcounter{utterance}  
    & & & \multicolumn{2}{p{0.3\linewidth}}{
        \cellcolor[rgb]{0.9,0.9,0.9}{
            \makecell[{{p{\linewidth}}}]{
                \texttt{\tiny{[GM$|$GM]}}
                {'word': 'ambivalence', 'assignment': 'team'}
            }
        }
    }
    & & \\ \\

    \theutterance \stepcounter{utterance}  
    & & & \multicolumn{2}{p{0.3\linewidth}}{
        \cellcolor[rgb]{0.9,0.9,0.9}{
            \makecell[{{p{\linewidth}}}]{
                \texttt{\tiny{[GM$|$GM]}}
                {'word': 'bonkers', 'assignment': 'team'}
            }
        }
    }
    & & \\ \\

    \theutterance \stepcounter{utterance}  
    & & & \multicolumn{2}{p{0.3\linewidth}}{
        \cellcolor[rgb]{0.9,0.9,0.9}{
            \makecell[{{p{\linewidth}}}]{
                \texttt{\tiny{[GM$|$GM]}}
                \texttt{opponent has won} \\
            }
        }
    }
    & & \\ \\

    \theutterance \stepcounter{utterance}  
    & \multicolumn{4}{p{0.6\linewidth}}{
        \cellcolor[rgb]{0.95,0.95,0.95}{
            \makecell[{{p{\linewidth}}}]{
                \texttt{\tiny{[P2$\rightarrow$P1]}}
                \texttt{Post{-}episode Reflexion (not sent to GM):} \\
\texttt{New plan: I failed because my clue "contradiction" was too abstract and directly matched the distractor "impulsivity" (or potentially "socialization" depending on the partner's interpretation), causing them to guess a distractor word and end their turn prematurely. The rule states that guessing a distractor ends the turn, so I must avoid any semantic overlap with the distractor list entirely. Next time, I will choose a concrete, action{-}oriented clue that is impossible to confuse with the distractors. For "bonkers," I will use "crazy" (if not a distractor) or "wild," and for "ambivalence," I will use "mixed" or "torn," ensuring these specific words do not appear in the distractor list (nationalization, impulsivity, socialization). I will also limit the targets to one per clue to maximize the chance of an immediate correct guess before the partner can accidentally hit a distractor.} \\
            }
        }
    }
    & & \\ \\

\end{supertabular}
}

\end{document}
